\chapter{从食物获得物质}\label{ch:g6:matter-from-food}

一头 \qty{40}{kg} 的小牛，两年后长成
\qty{600}{kg} 的母牛。半吨新增的牛体——肉、骨、皮、
角，从何而来？母牛吃的是草；新增的
物质\emph{并非草}；可是没有草，这些物质一样也不会存在。
\omterm{def:g1:animals-around-us:animal}{动物}如何处理食物，正是本章要讲的内容，
也就是
\cref{ch:g6:plants-are-producers}所记账目中\omterm{def:g5:food-webs:roles}{消费者}的那一面。

\section{动物用食物构建自身物质}

\begin{proposition}[消费者转化有机物]\label{prop:g6:matter-from-food:transform}
\omterm{def:g1:animals-around-us:animal}{动物}构建自身的\omterm{def:g6:plants-are-producers:matter}{有机物}——\omterm{def:g3:bones-and-muscles:muscle}{肌肉}、\omterm{def:g3:bones-and-muscles:skeleton}{骨骼}、
脂肪、\omterm{def:g1:animals-around-us:covering}{羽毛}、贝壳，原料取自\emph{所食的
\omterm{def:g6:plants-are-producers:matter}{有机物}}。\omterm{def:g4:journey-of-food:digestion}{消化}（\cref{def:g4:journey-of-food:digestion}）
把食物分解为\omterm{def:g4:journey-of-food:digestion}{营养物质}，\omterm{prop:g4:journey-of-food:absorb}{血液}将它们输送
（\cref{prop:g4:heart-and-blood:carries}），再由身体的
\omterm{def:g5:first-look-at-cells:cell}{细胞}（\cref{prop:g5:first-look-at-cells:growth}）重新组装，
变成\omterm{def:g1:animals-around-us:animal}{动物}自身的物质。\omterm{def:g1:animals-around-us:animal}{动物}是
物质的\emph{转化者}，绝不是从\omterm{def:g6:plants-are-producers:matter}{无机物}
中制造物质；这门本领只属于\omterm{def:g5:food-webs:roles}{生产者}。
\end{proposition}
\admitted

\begin{example}[草变成牛]\label{ex:g6:matter-from-food:cow}
跟随这半吨物质：母牛一吃就是几个小时
（\cref{prop:g2:what-animals-eat:daily}已说明为何这么久）；
她那繁重的\omterm{def:g4:journey-of-food:digestion}{消化}过程
（\cref{exo:g4:journey-of-food:11}）把草分解为\omterm{def:g4:journey-of-food:digestion}{营养物质}；
\omterm{prop:g4:journey-of-food:absorb}{血液}把它们运往乳房、\omterm{def:g3:bones-and-muscles:muscle}{肌肉}和
腹中正在生长的小牛；\omterm{def:g5:first-look-at-cells:cell}{细胞}将草的物质
转化为牛的物质。这片牧场，就字面而言，
正以新的形式四处走动。
\end{example}

\begin{remark}[物质易手，本质不变]\label{rem:g6:matter-from-food:hands}
在\omterm{def:g3:food-chains:chain}{食物链}的每一个箭头处，
（\cref{def:g3:food-chains:chain}），\omterm{def:g6:plants-are-producers:matter}{有机物}便易主，
并被重新构建——草的物质变为牛的物质，
牛的物质（\omterm{def:g2:animals-grow-up:born}{乳汁}、肉）变为儿童的物质。
\emph{\omterm{def:g6:plants-are-producers:matter}{有机物}}从不凭空产生：\omterm{def:g5:food-webs:roles}{消费者}将其
传递下去并加以改造。唯在\omterm{def:g3:food-chains:chain}{食物链}的绿色起点，它才首次被制造
（\cref{prop:g6:plants-are-producers:firstlink}）；
唯有在\omterm{def:g5:food-webs:roles}{分解者}的终点，即下一章的内容，它才
重新变为\omterm{def:g6:plants-are-producers:matter}{无机物}。
\end{remark}

\section{生长即物质生产}

\begin{definition}[动物的物质生产]\label{def:g6:matter-from-food:production}
\omterm{def:g1:animals-around-us:animal}{动物}的\emph{物质生产}\index{物质生产}，即其构建的新\omterm{def:g6:plants-are-producers:matter}{有机物}，包括幼体的
生长（小牛增重的千克数）、每年的更新（\omterm{def:g1:animals-around-us:covering}{皮毛}
脱换后再长出、鹿角、\omterm{def:g1:animals-around-us:covering}{羽毛}）、修复
（\cref{rem:g3:bones-and-muscles:alive}中愈合的\omterm{def:g3:bones-and-muscles:skeleton}{骨骼}），以及
产物（\omterm{def:g2:animals-grow-up:born}{乳汁}、卵、羊毛）。这一切都是由食物重新构建而来。
\end{definition}

\begin{example}[称量物质生产]\label{ex:g6:matter-from-food:weighing}
农民每天称量产量：一头奶牛每天产出
约 \qty{25}{kg} 的\omterm{def:g2:animals-grow-up:born}{乳汁}——都是她必须
日复一日从饲料中重新构建的\omterm{def:g6:plants-are-producers:matter}{有机物}；一枚鸡蛋
约含 \qty{60}{g} 的储备物质
（\cref{rem:g2:animals-grow-up:egg}），一夜之间由
她吃的谷物组装而成。无饲料，则无产出：账目
毫不留情地保持平衡。
\end{example}

\begin{proposition}[并非所有食物都变成身体]\label{prop:g6:matter-from-food:losses}
\omterm{def:g5:food-webs:roles}{消费者}从不会把全部食物转化为新物质：
\begin{itemize}
  \item 一部分从未被吸收，以粪便形式排出；
        （\cref{prop:g4:journey-of-food:absorb}）；
  \item 很大一部分被\emph{燃烧}——与
        \omterm{def:g4:breathing:movements}{呼吸}中的氧气结合，
        （\cref{prop:g4:breathing:exchange}）——为运动和保暖提供\omterm{prop:g3:balanced-diet:jobs}{能量}：
        物质作为燃料被消耗，
        以呼出的气体形式离开；
  \item 只有余下的部分才变为生长和产物。
\end{itemize}
母牛所吃的草中，只有其中一小部分变成新的
牛体——这正是为何
\cref{prop:g5:food-webs:pyramid}所示金字塔的每一层
都比其下面一层小得多。
\end{proposition}
\admitted

\begin{example}[金字塔终得解释]\label{ex:g6:matter-from-food:pyramid}
一片草地上数吨重的草，可养活一大群兔子；
兔子所含的物质，可养活一两只狐狸。
在每个箭头处，多数物质被燃烧用于生存，或被流失，
只有一小部分向上传递，成为身体。大型捕食者
（\cref{prop:g5:food-webs:pyramid}）之所以稀少，
是因为它们依靠\emph{剩余的剩余}为生——
这是算术，而非运气不佳。
\end{example}

\section{食谱、牙齿与同一门本领}

\begin{proposition}[食谱不同，转化相同]\label{prop:g6:matter-from-food:menus}
\omterm{def:g2:what-animals-eat:kinds}{植食性动物}、\omterm{def:g2:what-animals-eat:kinds}{肉食性动物}、\omterm{def:g2:what-animals-eat:kinds}{杂食性动物}
（\cref{def:g2:what-animals-eat:kinds}）指的是不同的
\omterm{def:g6:plants-are-producers:matter}{有机物}的\emph{来源}，而非不同的本领。
母牛用草重建自身，狐狸用兔子重建自身，二者都在执行
\cref{prop:g6:matter-from-food:transform}所述的转化
——其器官与食谱相匹配
（\cref{prop:g2:what-animals-eat:match}，
\cref{exo:g4:journey-of-food:11}）：草需要磨齿与庞大的\omterm{def:g4:journey-of-food:digestion}{消化}系统，
肉食则偏爱抓握器官与短肠。
\end{proposition}
\admitted

\begin{method}[核算动物的物质]\label{met:g6:matter-from-food:audit}
对任意\omterm{def:g1:animals-around-us:animal}{动物}，列出其物质账目：
\begin{enumerate}
  \item \emph{输入}：它摄取何种\omterm{def:g6:plants-are-producers:matter}{有机物}，又从
        谁那里获得？（即\omterm{def:g5:food-webs:web}{食物网}中指向它的箭头，
        见\cref{def:g5:food-webs:web}）；
  \item \emph{构建}：它生长、更新了哪些部分，
        又制造了哪些产物？
  \item \emph{输出}：粪便，以及为运动和保暖而
        燃烧掉的物质；
  \item 核对平衡：生产量必须只是
        很小的余量——若某一说法认为生产量更大（“这池塘
        产出的\omterm{prop:g3:sorting-living-things:groups}{鱼类}物质多于池塘所生长的食物”），
        则该说法有误。
\end{enumerate}
\end{method}

\begin{example}[核算的应用]\label{ex:g6:matter-from-food:audituse}
对乌鸫（见\cref{ex:g2:what-animals-eat:winter}）进行核算：
输入——蚯蚓、\omterm{prop:g3:sorting-living-things:groups}{昆虫}、浆果；构建——它自身的生长、
每年更换的新\omterm{def:g1:animals-around-us:covering}{羽毛}、一窝卵；输出——
墙上的粪便，以及大量被燃烧的部分，正是它们
为一年的飞行与鸣唱提供了\omterm{prop:g3:balanced-diet:jobs}{能量}。鸣禽每日的
食物重量相对体型极大，恰恰因为飞行
燃烧掉了账目中如此多的物质。
\end{example}

\section{练习}

\begin{exercise}[$\star$]\label{exo:g6:matter-from-food:1}
用一句话分别阐述\omterm{def:g5:food-webs:roles}{消费者}的生存方式，并将其与
\omterm{def:g5:food-webs:roles}{生产者}的生存方式进行对比。
\end{exercise}

\begin{exercise}[$\star$]\label{exo:g6:matter-from-food:2}
列出草转化为牛\omterm{def:g3:bones-and-muscles:muscle}{肌肉}的各个步骤，并指出
各步在哪一章中讲述。
\end{exercise}

\begin{exercise}[$\star$]\label{exo:g6:matter-from-food:3}
列举\omterm{def:g1:animals-around-us:animal}{动物}\omterm{def:g6:matter-from-food:production}{物质生产}的四种形式，每种各
举一例。
\end{exercise}

\begin{exercise}[$\star$]\label{exo:g6:matter-from-food:4}
说出\omterm{def:g5:food-webs:roles}{消费者}所摄取食物的三种去向，
依据\cref{prop:g6:matter-from-food:losses}。
\end{exercise}

\begin{exercise}[$\star$]\label{exo:g6:matter-from-food:5}
奶牛为何每天都需要饲料，才能持续
产\omterm{def:g2:animals-grow-up:born}{乳汁}？
\end{exercise}

\begin{exercise}[$\star$]\label{exo:g6:matter-from-food:6}
一枚卵中含有 \qty{60}{g} \omterm{def:g6:plants-are-producers:matter}{有机物}——它是
这些\omterm{def:g6:plants-are-producers:matter}{有机物}是在母鸡体内制成的——还是在
母鸡之前就已制成？请仔细厘清。
\end{exercise}

\begin{exercise}[$\star\star$]\label{exo:g6:matter-from-food:7}
运用\cref{prop:g6:matter-from-food:losses}，
解释食物金字塔为何逐层变窄。
\end{exercise}

\begin{exercise}[$\star\star$]\label{exo:g6:matter-from-food:8}
对一只宠物猫运用\cref{met:g6:matter-from-food:audit}进行核算：
输入、构建、输出——每项至少列出两条。
\end{exercise}

\begin{exercise}[$\star\star$]\label{exo:g6:matter-from-food:9}
“狐狸与奶牛对食物的处理方式各不相同。”
请运用
\cref{prop:g6:matter-from-food:menus}修正此说法——哪些不同，哪些
完全相同？
\end{exercise}

\begin{exercise}[$\star\star$]\label{exo:g6:matter-from-food:10}
一名儿童一年增重 \qty{3}{kg}，却吃掉了
数百千克食物。其余的去了哪里？请用
物质账目作答。
\end{exercise}

\begin{exercise}[$\star\star$]\label{exo:g6:matter-from-food:11}
蜗牛的壳属于\omterm{def:g6:plants-are-producers:matter}{无机物}质，然而
\cref{ex:g6:plants-are-producers:sorting}却将其算作
蜗牛自身的产物。请调和这一矛盾：蜗牛的身体
做了什么，借助何种原料才将其建成？
\end{exercise}

\begin{exercise}[$\star\star\star$]\label{exo:g6:matter-from-food:12}
一公顷土地种马铃薯所养活的人数，多于
用这些马铃薯喂猪后以猪肉为食所养活的人数。请运用
\cref{prop:g6:matter-from-food:losses}解释：马铃薯经过猪体绕行时，
这部分差额在马铃薯经由猪的
这一绕道中去了哪里？
\end{exercise}

\section{问题：农场账目}

\begin{problem}[{周末问题——对一座小农场进行核算，
从田间到餐桌}]\label{pb:g6:matter-from-food:1}
奶奶的小农场：一片牧场，养着一头奶牛，一座
以谷物饲养十只母鸡的鸡舍，一片菜地，以及
一家人的厨房。本周末，由你来记录这座农场的物质账目。

\medskip\noindent\textbf{第一部分——\omterm{def:g5:food-webs:roles}{生产者}层。}
\begin{enumerate}
  \item 列出该农场的\omterm{def:g5:food-webs:roles}{生产者}。它们的\omterm{prop:g4:what-plants-need:needs}{无机盐}来源是什么，
        \omterm{prop:g3:balanced-diet:jobs}{能量}来源又是什么？
  \item 牧场一季约生长 \qty{10000}{kg} 草的
        \omterm{def:g6:plants-are-producers:matter}{有机物}。这些草料中的每一克，
        是在谁的叶片中制造的？
  \item 奶奶把鸡舍粪便撒在菜地上。她遵守的是
        哪一条规则，取自
        \cref{prop:g5:people-and-nature:lasting}？
        粪便究竟归还了什么？
  \item 为何不能直接给奶牛喂\omterm{prop:g4:what-plants-need:needs}{无机盐}和水，
        使牧场无须参与？
\end{enumerate}

\medskip\noindent\textbf{第二部分——\omterm{def:g5:food-webs:roles}{消费者}层。}
\begin{enumerate}[resume]
  \item 奶牛每天大约吃掉 \qty{70}{kg} 草，
        产出约 \qty{20}{kg} 牛奶。请写出
        缺失的 \qty{50}{kg} 的三个去向。
  \item 母鸡把谷物转化为卵。写出母鸡的收支账目——
        输入、构建、输出——并说明正在换羽
        （全部\omterm{def:g1:animals-around-us:covering}{羽毛}更换）的母鸡为何
        在一段时间内产卵变少。
  \item 一只鹰叼走了一只雏鸡。将该农场的\omterm{def:g3:food-chains:chain}{食物链}延伸至鹰，
        并指出物质在哪些环节上被\emph{制造}，
        在哪些环节上只是转化。
  \item 这家人吃鸡蛋、牛奶和蔬菜。追溯一顿早餐——
        面包、水煮蛋、牛奶——中各项食物的\omterm{def:g6:plants-are-producers:matter}{有机物}，
        分别追溯到其最初被制造的
        叶片。
\end{enumerate}

\medskip\noindent\textbf{第三部分——各层级的
计算}
\begin{enumerate}[resume]
  \item 奶奶说：“牧场只能养活一头奶牛；若狐狸
        吃在此处以草饲养的兔子，它至多只能养活
        一个小型狐群。”请用各层级的损耗论证
        她的估算。
  \item 菜地直接养活这家人。运用\cref{exo:g6:matter-from-food:12}中的逻辑，
        解释为何菜地每平方米能养活的人数
        多于经由牧场、奶牛、牛奶这条途径
        所能养活的人数。
  \item 干旱时草停止生长。依次预测奶牛的牛奶、
        体重，以及收支账目中“构建”一行的变化，
        并说明降雨不足时，为何即使增加日照也
        无济于事
        （见\cref{pb:g6:plants-are-producers:1}中的第 7 题，
        其中给出了关键）。
  \item 每个层级用一句话结束收支账目：\omterm{def:g5:food-webs:roles}{生产者}、
        \omterm{def:g5:food-webs:roles}{消费者}——并指出账目中
        仍然缺少的是谁
        （即下一章的主题）。
\end{enumerate}
\end{problem}
